\usetikzlibrary{shapes.geometric,shapes.misc}

\begin{figure}[H]
\centering
\begin{tikzpicture}[
x=.5cm,
y=.5cm,
edge/.style={draw=black!55, line width=.25pt},
circlev/.style={circle, fill=black, draw=black, inner sep=1.4pt},
starv/.style={star, star points=5, star point ratio=2.1, fill=black, draw=black, inner sep=0pt, minimum size=6pt},
crossv/.style={cross out, draw=black, line width=.6pt, inner sep=0pt, minimum size=5pt},
squarev/.style={rectangle, fill=black, draw=black, inner sep=1.4pt},
squarev /.style={rectangle, fill=black, draw=black, inner sep=1.4pt},
]

\def\nodedata{
n1/1/16/circlev,
n2/3/13/circlev,
n3/3/18/circlev,
n4/4/14/circlev,
n5/5/20/circlev,
n6/7/17/circlev,
n7/6/11/starv,
n8/8/13/starv,
n9/10/15/starv,
n10/8/8/starv,
n11/9/9/starv,
n12/11/6/crossv,
n13/12/12/starv,
n14/13/8/crossv,
n15/15/10/crossv,
n16/13/3/crossv,
n17/14/4/crossv,
n18/16/1/squarev,
n19/17/7/crossv,
n20/18/3/squarev,
n21/20/5/squarev
}

\foreach[count=\i] \name/\xx/\yy/\style in \nodedata {
    \coordinate (\name) at (\xx,\yy);
    \expandafter\xdef\csname X\i\endcsname{\xx}
    \expandafter\xdef\csname Y\i\endcsname{\yy}
    \expandafter\xdef\csname Name\i\endcsname{\name}
}

\foreach \i in {1,...,20} {
    \pgfmathtruncatemacro{\jstart}{\i+1}
    \foreach \j in {\jstart,...,21} {
        \edef\Xi{\csname X\i\endcsname}
        \edef\Yi{\csname Y\i\endcsname}
        \edef\Xj{\csname X\j\endcsname}
        \edef\Yj{\csname Y\j\endcsname}
        \edef\NameI{\csname Name\i\endcsname}
        \edef\NameJ{\csname Name\j\endcsname}
        \pgfmathtruncatemacro{\cmp}{((\Xj>=\Xi)&&(\Yj>=\Yi)) || ((\Xi>=\Xj)&&(\Yi>=\Yj))}
        \ifnum\cmp=1
            \draw[edge] (\NameI) -- (\NameJ);
        \fi
    }
}

\foreach \name/\xx/\yy/\style in \nodedata {
    \node[\style] at (\name) {};
}

\draw[edge] (n1)  to[bend right=20] (n5);
\draw[edge] (n2)  to[bend right=20] (n6);
\draw[edge] (n7)  to[bend right=20] (n9);
\draw[edge] (n10) to[bend right=20] (n13);
\draw[edge] (n12) to[bend right=20] (n15);
\draw[edge] (n16) to[bend right=20] (n19);
\draw[edge] (n18) to[bend right=20] (n21);

\end{tikzpicture}
\caption{Three tiled copies of the width 3 graph gadget. The first copy excluding its right boundary is shown in $\bullet$. The second copy excluding its right boundary is shown in $\filledstar$. The third copy excluding its right boundary is shown in $\times$. The third copy's right boundary is shown in $\filledsquare$.}
\end{figure}